% p1-09-pellis-expansion

\section{P1 §9. Pellis Hierarchical Expansion}

9. Pellis Hierarchical Expansion

        Figure (fig-p1-ch9-pellis-expansion). Pellis expansion triptych --- additive sum of phi-inverse
                            powers / reciprocal interference / alpha-residue dial.

  9.1 Motivation
  The multiplicative grammar G$\varphi$ provides a zeroth-order symbolic representation but
  cannot represent values near lattice points without increasing C indefinitely. The Pellis
  hierarchical expansion introduces a nested additive correction scheme in powers of $\varphi$-1,
  providing finer-grained approximation without unbounded complexity growth. The
  construction draws on the structure of the Lucas ring Z[$\varphi$] (PhD: Ch.3 Trinity Identity)
  to keep partial sums algebraically tractable.

  9.2 Definition of the Expansion Class
  A hierarchical representation of a target quantity X $\in$ R^+ at depth K is defined as

  X approx F(K)($\varphi$) = sumk=1K ck $\varphi$-k, ck $\in$ Z, K $\in$ N. 34

  The formal asymptotic extension is

  F($\varphi$) $\sim$ sumk=1^$\in$fty ck $\varphi$-k as k $\rightarrow$ $\in$fty, 35

  interpreted as a formal series. The expansion variable $\varphi$-1 approx 0.6180 acts as a
  natural geometric-sequence parameter generating progressively finer corrections. The
  algebraic closure of Z[$\varphi$] ensures that integer-valued coefficients keep each partial sum
  in this ring.

  9.3 Two-Tier Structure
  The Pellis expansion is deployed as a second tier over the base multiplicative grammar:

  X approx F0 + Delta F(1) + Delta F(2) + …, Delta F(j) $\in$ span{$\varphi$-k}k $\geq$ 1, 36

  where F0 = n $\cdot$ 3k $\cdot$ $\varphi$p $\cdot$ pim $\cdot$ eq $\in$ S(C) is the zeroth-order best-fit from the base
  grammar and Delta F(j) are additive correction terms from a truncated Pellis expansion.

  9.4 Scale Ordering and Filtration
  The hierarchy induced by powers of $\varphi$-1 defines a natural scale ordering:

  $\varphi$-1 > $\varphi$-2 > $\varphi$-3 > …, 37

  with corresponding geometric ratio $\varphi$-1 approx 0.618. This induces a filtration

  F1 ⊃ F2 ⊃ F3 ⊃ …, Fk = spanR{$\varphi$-j}j $\geq$ k. 38

  9.5 Truncation Error and Convergence Criterion
  The truncation error at depth K is

  $\varepsilon$K(X) := | X - sumk=1K ck $\varphi$-k |. 39

  A hierarchical representation is operationally convergent if $\varepsilon$K(X) decreases
  monotonically with K and stabilizes below a predefined tolerance delta for some finite
  K0. No claim of analytic convergence of the infinite series is made; only finite-precision
  stabilization is asserted.

  9.6 MDL Penalty for the Hierarchical Expansion
  When the Pellis expansion is incorporated into the MDL framework, L(M1) is
  augmented by the cost of specifying expansion depth K and coefficient vector
  {ck}k=1K:

  L(M1Pellis) = L(M1base) + log2 K + sumk=1K log2(1 + |ck|). 40

  Only Pellis expansions with augmented MDL score below that of the null model are
  reported as compressive. This ensures that increased depth cannot manufacture
  spurious compression advantage.

